\documentclass[11pt]{article}
\usepackage{amsmath, amssymb, amsthm}

\title{Triangle Mutiny: Activation Algebra v2.7 (Generalized Foundations)}
\author{Borealis S. Hedling}
\date{17 October 2026}

\begin{document}
\maketitle

\section*{Overview}

This document provides the generalized algebraic foundations required for
Activation Formalism v2.8. It extends the completed algebra of v2.7 by adding
explicit domain declarations, operator typing, relation typing, bundle typing,
kernel hierarchy, interaction laws, semantic and provenance lifting schemas,
bundle composition, and the definition of the activation algebra object
$\mathfrak{A}$.

\section*{Algebraic Domains}

Define the following domains:

\[
\mathcal{M} = \text{Activation Models}
\]
\[
\mathcal{S} = \text{Symbolic Activation States}
\]
\[
\mathcal{B} = \text{Activation Bundles}
\]
\[
\mathcal{K} = \text{Activation Kernels}
\]

These domains form the basis of the activation algebra.

\section*{Operator Typing}

Define the activation operator:

\[
\mathcal{A} : \mathcal{M} \times \mathcal{S} \to \mathcal{M}
\]

\[
\mathcal{A} \in \text{End}(\mathcal{M})
\]

\section*{Relation Typing}

Define the activation relation:

\[
R : \mathcal{M} \times \mathcal{M} \to \mathbf{Bool}
\]

\[
R \subseteq \mathcal{M} \times \mathcal{M}
\]

\section*{Bundle Typing}

Define activation bundles:

\[
\mathcal{B} = \mathcal{A} \times R
\]

\[
\mathcal{B} \in \mathcal{A} \times R
\]

\section*{Kernel Hierarchy}

Define kernel hierarchy:

\[
K_{\text{geom}} \subseteq K_{\text{prov}} \subseteq K_{\text{sem}} \subseteq K_{\text{alg}}
\]

where:

\[
K_{\text{geom}}(M) = I_{\text{geom}}(M)
\]
\[
K_{\text{prov}}(M) = I_{\text{geom}}(M) \land I_{\text{prov}}(M)
\]
\[
K_{\text{sem}}(M) = K_{\text{prov}}(M) \land I_{\text{sem}}(M)
\]
\[
K_{\text{alg}}(M) = K_{\text{sem}}(M)
\]

\section*{General Interaction Laws}

\subsection*{Operator-Ordering Interaction}

\[
M_1 \preceq M_2 \Rightarrow
\mathcal{A}(M_1, S) \preceq \mathcal{A}(M_2, S)
\]

\subsection*{Relation-Operator Interaction}

\[
R(M_1, M_2) \Rightarrow
R(\mathcal{A}(M_1, S), \mathcal{A}(M_2, S))
\]

\subsection*{Kernel-Operator Interaction}

\[
K_{\text{alg}}(M) \Rightarrow K_{\text{alg}}(\mathcal{A}(M, S))
\]

\section*{Semantic Lifting Schema}

Define semantic lifting:

\[
\mathcal{L}_{\text{sem}} : \mathcal{M} \to \Sigma
\]

where $\Sigma$ is the semantic domain.

\[
\mathcal{L}_{\text{sem}}(\mathcal{A}(M, S)) = \mathcal{L}_{\text{sem}}(M)
\]

\section*{Provenance Lifting Schema}

Define provenance lifting:

\[
\mathcal{L}_{\text{prov}} : \mathcal{M} \to \Pi
\]

where $\Pi$ is the provenance domain.

\[
\mathcal{L}_{\text{prov}}(\mathcal{A}(M, S)) = \mathcal{L}_{\text{prov}}(M)
\]

\section*{Bundle Composition}

Define bundle composition:

\[
\mathcal{B}_1 \circ \mathcal{B}_2 =
(\mathcal{A}_1 \circ \mathcal{A}_2,\ R_1 \circ R_2)
\]

Composition is associative:

\[
(\mathcal{B}_1 \circ \mathcal{B}_2) \circ \mathcal{B}_3
=
\mathcal{B}_1 \circ (\mathcal{B}_2 \circ \mathcal{B}_3)
\]

\section*{Activation Algebra Object}

Define the activation algebra:

\[
\mathfrak{A} =
(\mathcal{M}, \mathcal{S}, \mathcal{A}, R, \preceq, \mathcal{B}, \mathcal{K},
 \mathcal{L}_{\text{sem}}, \mathcal{L}_{\text{prov}})
\]

This object is the foundation for the axioms introduced in v2.8.

\section*{Summary}

This generalized foundation provides:

\begin{itemize}
\item explicit algebraic domains,
\item operator, relation, and bundle typing,
\item kernel hierarchy,
\item general interaction laws,
\item semantic lifting schema,
\item provenance lifting schema,
\item bundle composition,
\item definition of the activation algebra object $\mathfrak{A}$.
\end{itemize}

These structures are required for the axiomatic activation formalism of v2.8.

\end{document}
